\documentclass{article}

\usepackage{amsmath}
\usepackage{algorithmic}
\usepackage{graphicx}
\usepackage{xspace}
\usepackage{multirow}
\usepackage{listings}
\usepackage{verbatim}
\begin{document}
\title{Hello \LaTeX\xspace world}

\author{Sergio Vasquez\\
\small{Computer Engineering and Computer Science}\\
\small{\texttt{sergio.vasquez01@student.csulb.edu}}}
\date{}
\maketitle

\begin{abstract}
This document is a model and instructions for \LaTeX\xspace `article' class
\end{abstract}

\section{Introduction}
Welcome to the \LaTeX\xspace world.
\section{Ease of Use}

\subsection{Maintaining the Integrity of the Specifications}
The `article' class is used to format your paper and style the text. All margins, column widths, line spaces, and text fonts are prescribed.

\section{Styling Guide}

\subsection{Abbreviations and Acronyms}
Define abbreviations and acronyms the first time they are used in the text, 
even after they have been defined in the abstract.

\subsection{Equations}
\begin{equation}
\sum_{n=0}^{\infty}\frac{af^{n}}{n!}(x-a)^n
\label{taylor}
\end{equation}
\eqref{taylor} is the famous Taylor series. Use ``\eqref{taylor}'', not ``Eq.~\eqref{taylor}'' or ``equation \eqref{taylor}'', except at the beginning of a sentence: ``Equation \eqref{taylor} is . . .''

Taylor series in a text would be $\sum_{n=0}^{\infty}\frac{af^{n}}{n!}(x-a)^n$.

\subsection{Lists}
Bullet style list.
\begin{itemize}
    \item item 1
    \item item 2
    \item item 3
\end{itemize}
Number style list.
\begin{enumerate}
    \item item 1
    \item item 2
    \item item 3
\end{enumerate}

\subsection{Figures and Tables}
\paragraph{Positioning Figures and Tables} Figure captions should be below the figures; table heads should appear above the tables. Insert figures and tables after they are cited in the text. Use the abbreviation ``Fig. \ref{fig}''.

\begin{table}[htpb]
\caption{Table Type Styles}
\begin{center}
\begin{tabular}{|c||c|c|c|}
\hline
\multirow{2}{*}{\textbf{\begin{tabular}[c]{@{}c@{}}Table\\ Head\end{tabular}}} & \multicolumn{3}{c|}{\textbf{Table Column Head}}                                                \\ \cline{2-4} 
                                                                               & \textit{\textbf{Table column subhead}} & \textit{\textbf{Subhead}} & \textit{\textbf{Subhead}} \\ \hline
\multicolumn{1}{|l||}{}                                                         &                                        &                           &                           \\ \hline
\end{tabular}
\end{center}
\end{table}

\begin{figure}[htbp]
\centering
\includegraphics[width=0.5\columnwidth]{fig1.png}
\caption{Working example}
\label{fig}
\end{figure}

\subsection{Algorithms}
\begin{algorithmic}
    \STATE $i\gets 10$
    \IF {$ i \geq 5$}
    \STATE $i\gets i -1$
    \ELSE 
    \IF {$i \leq 3$}
    \STATE $i \gets i + 2$
    \ENDIF
    \ENDIF
\end{algorithmic}

\subsection{Source codes}
\noindent\rule{\textwidth}{0.4pt}
\vspace{-5ex}
\verbatiminput{HelloWorld.java}
\vspace{-3ex}
\noindent\rule{\textwidth}{0.4pt}
\subsection{References}
Please number citations consecutively within brackets
\cite{Eason}. The sentence punctuation follows the bracket
\cite{Maxwell}. Refer simply to the reference number, as in \cite{Jacobs}---do not use ``Ref. \cite{Jacobs}'' or ``reference \cite{Jacobs}'' except at the beginning of a sentence.

\begin{thebibliography}{00}
\bibitem{Eason} G. Eason, B. Noble, and I. N. Sneddon, ``On certain integrals of Lipschitz-Hankel type involving products of Bessel functions,'' Phil. Trans. Roy. Soc. London, vol. A247, pp. 529--551, April 1955.
\bibitem{Maxwell} J. Clerk Maxwell, A Treatise on Electricity and Magnetism, 3rd ed., vol. 2. Oxford: Clarendon, 1892, pp.68--73.
\bibitem{Jacobs} I. S. Jacobs and C. P. Bean, ``Fine particles, thin films and exchange anisotropy,'' in Magnetism, vol. III, G. T. Rado and H. Suhl, Eds. New York: Academic, 1963, pp. 271--350
\end{thebibliography}

\end{document}
